\documentclass[11pt]{article}
\usepackage[margin=1in]{geometry}
\usepackage{amsmath,amssymb,amsthm,hyperref,enumitem}
\newtheorem{theorem}{Theorem}
\newtheorem{lemma}{Lemma}
\title{The first positive Mittag--Leffler zero near order one}
\author{Principia Math In Progress}
\date{21 September 2026}
\begin{document}
\maketitle
\noindent\textbf{Task 1: MathDB 361626.}
For $a>0$, put
\[
 f_a(t)=E_{a,1}(-t^a)=\sum_{k=0}^{\infty}
             \frac{(-1)^k t^{ak}}{\Gamma(ak+1)},\qquad t\geq0,
\]
and let $Z_1(a)$ be its smallest positive zero, when one exists.
Conjecture 3(b) of Chaudhary, Diethelm and Hashemishahraki
\cite{CDH} proposes, as $a\downarrow1$,
\[
 Z_1(a)=(c\log(a-1)+d)(1+o(1)),
 \qquad c\approx-0.275,\quad d\approx17.54.
\]
The following theorem disproves the proposed leading coefficient and
gives the asymptotic with an additive error tending to zero.

\begin{theorem}\label{main}
For all sufficiently small $\varepsilon>0$, $f_{1+\varepsilon}$ has a
smallest positive zero. With $L=\log(1/\varepsilon)$,
\[
 \boxed{Z_1(1+\varepsilon)=L+\log L+o(1).}
\]
In particular $Z_1(1+\varepsilon)/L\to1$.
\end{theorem}

The conjectured leading coefficient is therefore forced to be $c=-1$,
far from its stated numerical fit. This is a disproof of part (b), and
hence of the conjunction recorded in MathDB. The global positivity
bound in part (a) and the endpoint expansion at $a=2$ in part (c) are
not settled here. A constant $d$ is invisible at the conjecture's
multiplicative $1+o(1)$ precision; no contradiction is based on $d$.

\section*{A classical integral representation}
Throughout the proof $a=1+\varepsilon$, $0<\varepsilon<1/4$, and
$\theta=\pi\varepsilon$. Define the positive density
\[
 K_\varepsilon(r)=\frac{\sin\theta}{\pi}
 \frac{r^{a-1}}{(r^a-1)^2+2r^a(1-\cos\theta)},\qquad r>0.
\]
The standard residue and branch-cut representation is
\begin{equation}\label{representation}
 f_a(t)=\frac2a e^{t\cos(\pi/a)}\cos(t\sin(\pi/a))
             -\int_0^\infty e^{-tr}K_\varepsilon(r)\,dr
 \quad(t>0).
\end{equation}
For completeness, termwise Laplace transformation of the defining
series for real $s>1$ gives
\[
 \int_0^\infty e^{-st}f_a(t)\,dt
 =\sum_{k\geq0}(-1)^k s^{-ak-1}
 =\frac{s^{a-1}}{s^a+1}.
\]
Use the principal branch, cut along the negative real axis. There are
two simple poles at $e^{\pm i\pi/a}$; their inverse-Laplace residues
are $a^{-1}\exp(t e^{\pm i\pi/a})$. On the upper side of the cut,
\[
 \frac{s^{a-1}}{s^a+1}\bigg|_{s=-r+i0}
 =\frac{r^{a-1}e^{i\theta}}{1-r^ae^{i\theta}},\qquad
 \frac1\pi\operatorname{Im}\frac{r^{a-1}e^{i\theta}}
                                     {1-r^ae^{i\theta}}
 =K_\varepsilon(r).
\]
Deforming the inversion contour to the cut consequently gives the
negative integral in \eqref{representation}. The small-circle
contribution is $O(\rho^a)$ as its radius $\rho\downarrow0$;
the large arcs vanish for $t>0$ since the transform is $O(1/s)$
and the exponential decays on the left arcs. The cut integral is
absolutely convergent at both endpoints.

The exact total mass, obtained with $u=r^a$, is
\begin{equation}\label{mass}
 \int_0^\infty K_\varepsilon(r)\,dr
 =\frac{\sin\theta}{\pi a}\int_0^\infty
      \frac{du}{(u-\cos\theta)^2+\sin^2\theta}
 =\frac{\pi-\theta}{\pi a}=\frac2a-1.
\end{equation}
The antiderivative used here is
$(\pi a)^{-1}\arctan((u-\cos\theta)/\sin\theta)$.
The representation and the exponential/algebraic decomposition are
classical; see, for example, \cite{Hanneken,Paris02,Paris21}.
Our argument below supplies the needed uniform bounds and identifies
the \emph{first} zero in the simultaneous limit.

\section*{Uniform expansion in a logarithmic window}
\begin{lemma}\label{window}
Uniformly for $L/4\leq t\leq3L/2$, as $\varepsilon\downarrow0$,
\begin{equation}\label{uniform}
 f_a(t)=e^{-t}(1+o(1))
             -\varepsilon\Gamma(a)t^{-a}(1+o(1)).
\end{equation}
The two errors in this formula are separately uniform relative errors.
\end{lemma}
\begin{proof}
All constants in the following estimates can be chosen independently
of $\varepsilon$ and $t$ in this window. Put $h=t^{-2}$, so eventually
$h<1/2$. Write $D(r)=1-2r^a\cos\theta+r^{2a}$.
For $0\leq r\leq1/2$, $D(r)\geq(1-r^a)^2\geq1/4$ and
$D(r)^{-1}=1+O(r^a)$. Thus
\begin{align*}
 \int_0^{1/2}e^{-tr}K_\varepsilon(r)\,dr
 &=\frac{\sin\theta}{\pi}
   \left\{\Gamma(a)t^{-a}
       +O(t^{-2a})+O(t^{-a}(1+t)e^{-t/2})\right\}\\
 &=\varepsilon\Gamma(a)t^{-a}(1+o(1)).
\end{align*}
Here the $O(t^{-2a})$ bound follows by extending the error integral
to infinity; its factor $\Gamma(2a)$ is uniformly bounded. For the
main gamma integral the omitted tail is bounded, after $u=tr$, by
$t^{-a}\int_{t/2}^{\infty}(1+u)e^{-u}\,du$.

Next, the unweighted $K_\varepsilon$ mass outside $[1/2,3/2]$ is
$O(\varepsilon)$. The lower part follows from $D\geq1/4$; on the
upper part $D\geq(r^a-1)^2\geq r^{2a}/9$, and integration of
$O(\varepsilon r^{-a-1})$ gives the assertion. On $[1/2,3/2]$,
$r^{a-1}$ is bounded and the mean value theorem gives
$|r^a-1|\geq C_0|r-1|$ with a fixed $C_0>0$. Consequently
\[
 K_\varepsilon(r)\leq
            \frac{C\varepsilon}{(r-1)^2}\quad(r\ne1).
\]
Combining this estimate with \eqref{mass} yields
\[
 \int_{1-h}^{1+h}K_\varepsilon(r)\,dr
 =1+O(\varepsilon)+O(\varepsilon/h)=1+o(1),
\]
since $\varepsilon t^2\to0$ uniformly. On this central interval
$e^{-tr}=e^{-t}(1+O(th))$; hence its weighted contribution is
\[
 \int_{1-h}^{1+h}e^{-tr}K_\varepsilon(r)\,dr
 =e^{-t}(1+o(1)).
\]
On the rest of $[1/2,3/2]$ the contribution is at most
$C\varepsilon h^{-2}e^{-t/2}=C\varepsilon t^4e^{-t/2}$.
Dividing by $e^{-t}$ bounds it by
\[
 C\varepsilon(3L/2)^4 e^{3L/4}=O(L^4\varepsilon^{1/4})=o(1).
\]
The remaining tail $r\geq3/2$ contributes at most
$C\varepsilon e^{-3t/2}=o(e^{-t})$.

Finally, $\cos(\pi/a)=-1+O(\varepsilon^2)$ and
$\sin(\pi/a)=O(\varepsilon)$. Since $\varepsilon t\to0$,
the residue term in \eqref{representation} is $2e^{-t}(1+o(1))$.
Subtracting the central contribution $e^{-t}(1+o(1))$ and the
small-$r$ contribution proves \eqref{uniform}.
\end{proof}

\section*{Excluding earlier zeros}
\begin{lemma}\label{early}
For all sufficiently small $\varepsilon>0$, $f_a(t)>0$ for
$0\leq t\leq L/4$.
\end{lemma}
\begin{proof}
We compare directly with $f_1(t)=e^{-t}$. For $1\leq b\leq a$
and $k\geq1$,
\[
 \frac{\partial}{\partial b}
 \frac{t^{bk}}{\Gamma(bk+1)}
 =\frac{k t^{bk}}{\Gamma(bk+1)}
                         (\log t-\psi(bk+1)).
\]
The standard gamma bounds $\Gamma(bk+1)\geq k!$ and
$0<\psi(bk+1)\leq\log(bk+1)\leq2k$ apply. The first uses
monotonicity of $\Gamma$ on $[2,\infty)$; the digamma upper bound
also follows from Jensen's inequality for the logarithm under a
gamma density. These estimates justify differentiating the series
on compact parameter and $t$ intervals.

For $1\leq t\leq L/4$, eventually $t^\varepsilon\leq2$.
Put $T=t^b\leq2t$. Summing absolute values of the derivatives gives
\begin{align*}
 |\partial_b f_b(t)|
 &\leq\sum_{k\geq1}\frac{(k\log t+2k^2)T^k}{k!}\\
 &=(T\log t+2T^2+2T)e^T
 \leq C(1+t)^3e^{2t}.
\end{align*}
The fundamental theorem of calculus therefore implies
\[
 \frac{|f_a(t)-e^{-t}|}{e^{-t}}
 \leq C\varepsilon(1+L)^3e^{3L/4}
 =O(\varepsilon^{1/4}(1+L)^3)=o(1).
\]
For $0<t\leq1$, use
$k t^{bk}|\log t|\leq1/(b\mathrm e)\leq1/\mathrm e$.
The derivative series is then uniformly bounded by
$\sum_{k\geq1}(1/\mathrm e+2k^2)/k!$, giving
$|f_a(t)-e^{-t}|=O(\varepsilon)$ uniformly. Since $e^{-t}\geq e^{-1}$
there and $f_a(0)=1$, positivity follows on the entire interval.
\end{proof}

\section*{Locating the first zero}
Fix $0<\delta<1/2$. On $[L/4,(1-\delta)L]$ consider the ratio
of the second leading term in \eqref{uniform} to the first:
\[
 R_\varepsilon(t)=\varepsilon\Gamma(a)e^t t^{-a}.
\]
For large $L$ this is increasing throughout the interval because
its logarithmic derivative is $1-a/t>0$. At its upper endpoint
it is $O(e^{-\delta L}L^{-a})=o(1)$. Lemma~\ref{window}
therefore gives positivity throughout this interval. Together with
Lemma~\ref{early} this rules out every zero up to $(1-\delta)L$.
At $t=(1+\delta)L$, the ratio is
$\Gamma(a)e^{\delta L}((1+\delta)L)^{-a}\to\infty$, so
\eqref{uniform} gives a negative value. Continuity now supplies a
zero, and the zero set has a least positive element since it is
closed and bounded away from zero. It follows that
\[
 (1-\delta)L<Z_1(a)<(1+\delta)L.
\]
As $\delta$ was arbitrary, $Z_1(a)/L\to1$.

At this first zero, both relative errors in \eqref{uniform} tend
to zero. Taking logarithms of the equality of the two terms gives
\[
 Z_1(a)=L+a\log Z_1(a)-\log\Gamma(a)+o(1).
\]
Now $\log Z_1(a)=\log L+o(1)$,
$\varepsilon\log Z_1(a)=o(1)$, and $\log\Gamma(a)=o(1)$.
This proves Theorem~\ref{main}.
\hfill$\square$

\section*{Verification and scope}
The proof is analytic and does not infer uniformity or first-zero
minimality from numerical samples. The accompanying verifier checks
the branch-cut algebra and mass antiderivative symbolically. It also
performs supplemental 100-digit power-series calculations for five
values of $\varepsilon$ from $10^{-2}$ to $10^{-20}$, checks the
predicted sign transition, and compares two values against independent
quadrature of \eqref{representation}. These computations are not
interval certificates. Review is by the authoring agent only.
The complete current preprint and the complete author-posted text
of its journal final were read; the bounded priority audit is
documented separately and does not certify exhaustive novelty.

\section*{Four proposed follow-up tasks}
\begin{enumerate}[label=\textbf{Task \arabic*:},start=2,leftmargin=*]
\item Determine the next term after $L+\log L$, including whether
it is $(\log L-2)/L+o(1/L)$.
\item Give explicit constants and an explicit $\varepsilon_0$
in a two-sided error bound for the first-zero asymptotic.
\item Prove or refute the source's global bound
$f_a(t)>0$ for $1<a\leq2$ and $0\leq t\leq1.559$.
\item Derive a closed expression for the coefficient of $2-a$
in $Z_1(a)$ at $a=2$ and rigorously compare it with the source's fit.
\end{enumerate}

\small
\begin{thebibliography}{9}
\bibitem{CDH} R. Chaudhary, K. Diethelm and S. Hashemishahraki,
On the separation of solutions to fractional differential equations
of order $\alpha\in(1,2)$, \emph{Applied Numerical Mathematics}
\textbf{203} (2024), 84--96,
\href{https://doi.org/10.1016/j.apnum.2024.05.020}{doi:10.1016/j.apnum.2024.05.020}.
Current preprint: \href{https://arxiv.org/abs/2401.14771}{arXiv:2401.14771v3}.
Conjecture 3(b), journal p.~95; preprint p.~17.
\bibitem{Hanneken} J. W. Hanneken, B. N. N. Achar and D. M. Vaught,
An alpha-beta phase diagram representation of the zeros and properties
of the Mittag-Leffler function, \emph{Advances in Mathematical Physics}
(2013), Article 421685,
\href{https://doi.org/10.1155/2013/421685}{doi:10.1155/2013/421685}.
\bibitem{Paris02} R. B. Paris, Exponential asymptotics of the
Mittag--Leffler function, \emph{Proceedings of the Royal Society A}
\textbf{458} (2002), 3041--3052,
\href{https://doi.org/10.1098/rspa.2002.0977}{doi:10.1098/rspa.2002.0977}.
\bibitem{Paris21} R. B. Paris, Asymptotics of the Mittag--Leffler
function $E_a(z)$ on the negative real axis when $a\to1$,
\href{https://arxiv.org/abs/2005.05737}{arXiv:2005.05737v2} (2021).
This paper restricts $a<1$; its first-zero counterpart is not asserted.
\end{thebibliography}
\end{document}
